
\textbf{Encoder \& Decoder Architecture}~
The encoder is constructed as a sequential series of components: starting with a 1D convolutional layer featuring $C$ channels and a kernel size of 7, followed by a set of $B$ residual conventional blocks. Each block is composed of two dilated convolutions with dilation rate of $(1,1)$ and kernel size of $(3,1)$ and a skip-connection, followed by a strided convolutional down-sampling layer, with a kernel size $K$ of the twice the stride $R$. Whenever down-sampling, the number of channels is doubled. Unlike in EnCodec that the convolution blocks are followed by a two-layer LSTM, we use BiLSTM to augment the semantic modeling ability. A final 1D convolution layer with a kernel size of 7 is used to set the dimensionality of embeddings to $D$. We use $C=32,B=4$ and $(2,4,5,8)$ as strides. We use ELU~\citep{clevert2016fast} as a non-linear activation either layer normalization~\citep{ba2016layer} or weight normalization~\citep{salimans2016weight}.The decoder mirrors the encoder and uses transposed convolutions and LSTM instead of stride convolutions and BiLSTM, with the strides in reverse order as in the encoder. The decoder outputs the final audio signal.

\textbf{Residual Vector Quantizer}~
We use Residual Vector Quantizer~(RVQ) to quantize the encoder output and follow the same training procedure as EnCodec. During training, the code selected for each input is updated using an exponential moving average with a decay of 0.99, and codes which have not been assigned any input vector for several batches are replaced with input vectors randomly sampled within current batch. Straight-through-estimator~\citep{bengio2013estimating} is used to compute the gradient of encoder, e.g. as if the quantization step was the identity function during the backward phase. Finally, a commitment loss, consisting of the MSE between the input of the quantizer and its output, with gradient only computed with respect to its input, is added to the overall training loss.

\textbf{Discriminator}~
The MS-STFT discriminator utilizes networks with identical structures that operate on multi-scaled complex-valued STFT, where the real and imaginary parts are concatenated. For each sub-network, it is composed of a 2D convolutional layer (using kernel size $3\times8$ with 32 channels), followed by 2D convolutions with increasing dilation rates in the time dimension (1, 2 and 4), and a stride of 2 over the frequency axis. A final 2D convolution with kernel size $3\times3$ and stride $(1, 1)$ provide the final prediction. For MSD and MPD, we follow the same settings as in HiFiGAN~\citep{kong2020hifi} but adjust the channel number to align the discriminator's parameters more closely with that of MS-STFT.